% 方法流程图的共用样式：黑白矩形框、直角箭头、宋体正文、侧边竖排分区标题。
% 坐标单位为毫米，y 轴向下；版心宽 158 mm，侧边标题占 0--6 mm，正文框占 9--158 mm。
\usetikzlibrary{arrows.meta,calc}
% ctex 默认行距系数为 1.3，图内统一取 1。
\newcommand{\fcsize}[2]{\linespread{1}\fontsize{#1}{#2}\selectfont}
\tikzset{
  fc/.style={x=1mm,y=-1mm,>=Latex,line width=0.45pt},
  box/.style={draw,rectangle,anchor=north west,inner sep=1.5mm,outer sep=0pt,align=left,
    font=\songti\fcsize{8.6}{11}},
  % 固定正文高度并顶端对齐，使同一行各框的标题对齐；框外高为 h + 3 mm。
  h/.style={execute at begin node={\begin{minipage}[t][#1][t]{\linewidth}\raggedright},
    execute at end node={\end{minipage}}},
  io/.style={box,dashed},
  key/.style={box,line width=0.9pt},
  edge/.style={->,line width=0.45pt},
  note/.style={font=\songti\fcsize{8}{9.5},fill=white,inner sep=0.6mm},
}
% 框标题
\newcommand{\bt}[1]{{\heiti #1}\\}
% 侧边分区：\sidelabel{上沿 y}{下沿 y}{竖排文字}
\newcommand{\sidelabel}[3]{%
  \draw (0.3,#1) rectangle (6.1,#2);
  \node[align=center,font=\heiti\fcsize{9.5}{11.5},inner sep=0pt]
    at (3.2,{(#1+#2)/2}) {#3};}
% 总线连接：上方节点列表、下方节点列表、总线 y、总线左端 x、右端 x
\newcommand{\bus}[5]{%
  \foreach \u in {#1} {\draw (\u.south) -- (\u.south |- 0,#3);}
  \draw (#4,#3) -- (#5,#3);
  \foreach \l in {#2} {\draw[edge] (\l.north |- 0,#3) -- (\l.north);}}
